% Compact main-article section.  Do not also input port_boundary_theorems.tex.
% Expected bibliography keys: AHT2006, Lackenby2021, Lackenby2026.
\section{Ordered interfaces and safe disc-bundle replacement}
\label{sec:parallelity-interface}

The maintained source parallelity packet already describes a cut as a
bounded inherited core together with interval bundles whose bases are
binary families of polygons.  It retains all selected 0-, 1- and 2-handles,
including disc bases supplied only by a selected 2-handle.  Its component
profiles record Euler characteristic, orientability, boundary count and
core-facing ports, but do not extract the ports' cyclic boundary order.
Walking the subdivided boundary can take time proportional to normal weight.
We supply an executable extraction of the finite interface and distinguish
the correct replacements for fully and partially attached disc bundles.

These are additions to this implementation, not claims of a new general
parallelity algorithm.  Lackenby's Theorems 9.2 and 9.3 already provide
polynomial algorithms with attachment data, including cyclic order around
vertical annuli \cite{Lackenby2021}.  Report 68 in the repository already
proves the degree-two ordered-boundary principle, recovering cut-path
endpoints from their count, sum, square sum and path length.  Our contribution
is its integration with the later all-dimensional source packet, endpoint
occurrences at coincident vertices, zero-gap flags, fibre transport, correct
partial-disc replacement, and independently replayable code.

\subsection{Extracting the port matching}

Let $B$ be a finite compact surface represented by $s$ interval records of
polygon copies and affine side gluings.  There are $N$ corner addresses.
Let $V$ denote the corner-pairing system induced by glued polygon sides;
its orbits are quotient vertices.  Distinguish $p$ individual free sides,
the \emph{ports}, each specified by a binary rank.  Port $j$ has two endpoint
occurrences $e_{j,0},e_{j,1}$.  These labels remain distinct when their
addresses represent the same vertex.  Other free sides are old-boundary
sides.  Let $R$ be $V$ together with the endpoint pairings of old-boundary
sides, omitting the endpoint pairings of ports.

\begin{theorem}[Finite interface extraction]
\label{thm:finite-interface}
Each $R$-orbit containing marked endpoints contains exactly two endpoint
occurrences.  Their pairing $\tau$, and the information whether each paired
gap has zero length, determine all cyclic port words.  They are computable
using $O(p^2)$ ordinary interval-orbit counting calls on systems with
$O(s+p)$ records and universe size $N$.  Hence an algorithm exists whose
running time is polynomial in $s+p+\log(N+1)$; the output has $O(p)$
combinatorial records.
\end{theorem}
\begin{proof}
After quotienting by $V$, the free sides cellulate the disjoint union of
circles $\partial B$.  Delete the open port sides.  Each remaining component
is a circle with no marked endpoint, or an interval with two marked endpoint
occurrences.  A degenerate interval may be a single vertex, as when ports
are adjacent or a single loop edge is deleted.  The $R$-orbits describe
exactly these components.  Thus the marked orbits give a perfect matching
$\tau$ of the endpoint occurrences.

For any pairing system $Q$, let $c(Q)$ be its orbit count.  Adding one
singleton pairing gives
\[
 c(Q\cup\{x\leftrightarrow y\})=
 \begin{cases}c(Q),&x\sim_Qy,\\c(Q)-1,&x\not\sim_Qy.\end{cases}
\]
One edge merges precisely two previously distinct components, or none.
Testing all pairs of marked addresses therefore recovers $\tau$ using
$O(p^2)$ calls.  For a matched pair, the analogous test in $V$ detects a
zero-length gap: its ends represent the same vertex exactly when the
remaining interval contains no old-boundary side.  A nondegenerate interval
cannot have its ends identified in a cellulation of a one-manifold.

Let $\sigma$ exchange the two occurrences belonging to each port.
The multigraph with the $\sigma$ and $\tau$ pairs as edges is 2-regular.
Its alternating cycles are exactly the boundary circles containing ports.
Traversing them gives the cyclic words in $O(p)$ time after sorting labels;
choose the least port label and its endpoint zero to fix each direction.
The polynomial bound follows from the AHT orbit-counting theorem
\cite{AHT2006}.  When $p=0$ the answer is empty and no counting call is needed.
\end{proof}

For the maintained source packet, $s=O(T)$, $p\leq66T$, and
$\log(N+1)=O(L+\log(T+1))$, where $T$ is the tetrahedron count and $L$
bounds the coordinate bit length.  The theorem therefore yields a
polynomial reduction in encoded source size.  This does not identify the
repository's particular interval implementation with the proved AHT bound:
its currently established replay bound is in the realized trace size.

The prototype uses two marked weighted queries in place of repeated orbit
counts.  Endpoint $e_{j,a}$ receives weight $2^{2j+a}$ at its corner address;
overlapping weights are added.  Positive weights of $R$-orbits have exactly
two set bits, and positive weights of $V$-orbits have one or two.  Each
endpoint occurs once, so no bit carries arise and every positive mask has
multiplicity one.  All unmarked orbits share the zero-weight row.  The
weights have $2p$ bits, and both positive histograms have $O(p)$ rows.

\subsection{Fibre transport and the attachment subset}

Each polygon gluing has a fibre-reversal bit $\epsilon$.  Lift its corner
map to two sheets by
\[
                    (x,a)\longmapsto(g(x),a\oplus\epsilon).
\]
Old-boundary side maps preserve the fibre in their polygon charts.
Applying the matching construction to the lifted system gives the fibre
change along each gap, by following sheet zero at one end to its sheet at
the other.  This is the prototype's third weighted query, using $4p$
endpoint bits.

\begin{proposition}[Signed consistency and attachment components]
\label{prop:interface-consistency}
The lifted matching commutes with the deck involution.  For an orientable
total interval bundle, the xor of gap fibre changes around every marked
boundary circle is zero.  Let $A\subseteq\partial B$ be the union of port
sides.  A recovered circle lies entirely in $A$ exactly when every gap has
zero length.  On any other recovered circle, the number of components of
$A$ is the number of nonzero gaps.
\end{proposition}
\begin{proof}
An interval bundle over an interval or a point is trivial, so each gap has
two distinct lifts exchanged by the deck involution.  In an orientable
total space, the fibre orientation character equals that of the base.
The restriction of the tangent bundle of the base to its boundary splits
as the tangent line of a circle and an outward normal line, both orientable.
Thus fibre holonomy is zero on every boundary circle, including boundaries
of nonorientable bases.  Port traversal itself preserves the local fibre,
giving the xor assertion.  Finally, delete the nonzero gaps in a circular
word: either none is deleted and $A$ is the full circle, or each such gap
separates one of its interval components from the next.
\end{proof}

This does not imply that the full bundle is a product: an orientable
interval bundle over a M\"obius band has nontrivial fibre holonomy on its
core curve and trivial holonomy on its boundary.  Nor do independently
directed cyclic words supply coherent orientations between different
boundary circles or arbitrary markings on horizontal boundaries.  Those
data must be carried separately when a patterned hierarchy requires them.
If a base profile has $b$ boundary circles and $c$ recovered marked circles,
its unmarked circles remain the binary count $b-c$.

\subsection{Disc products with full and partial attachments}

Let a disc product $E=D^2\times I$ meet the complementary core $K$ in
$A\times I$, with its source attachment charts, where $A$ is a union of
port sides in $\partial D^2$.  The interfaces are embedded, and interfaces
of different bundle components are disjoint.  These are hypotheses supplied
by the authenticated source packet.

\begin{theorem}[Safe local replacement]
\label{thm:partial-disc-replacement}
If $A=\varnothing$, $E$ is an isolated 3-ball.  If $A=\partial D^2$, adjoining
$E$ attaches one 2-handle.  Otherwise let $r\geq1$ be the number of interval
components of $A$.  Adjoining $E$ is, after transport through collars of the
attaching discs, the addition of $r-1$ relative 1-handles.  The homeomorphism
can be fixed outside those collars and transports all explicitly retained
boundary markings.  Over all attached disc bases, the number of relative
handles added is at most $p$.
\end{theorem}
\begin{proof}
The empty and full-boundary cases follow directly from the product
description.  In the remaining case, choose $r-1$ disjoint properly embedded
arcs in $D^2$ with endpoints outside $A$, cutting it into $r$ discs, each
containing precisely one interval of $A$ on its boundary.  Such arcs are
constructed by successively cutting off the intervals in cyclic order;
their dual graph is a tree.  Their products with $I$ are vertical discs
splitting $E$ into $r$ balls.  Each ball meets $K$ in one disc and can be
absorbed into a collar of that disc.  The enlarged core is homeomorphic to
$K$ by a homeomorphism fixed away from the attaching collars.  Rejoining
across the $r-1$ cut discs amounts to attaching $r-1$ 1-handles between
their collar patches.  The same charts transport boundary markings.

A full attachment uses at least one port and contributes one handle.
A partial attachment using $p_i$ ports has $r\leq p_i$ and contributes
$r-1\leq p_i$ handles.  Sum over components.  Unattached balls can occur
with large multiplicity and must be counted rather than listed.
\end{proof}

The Euler contribution is $1-\chi(A)$: it is $1$ for an empty or full
attachment, and $1-r$ for a proper nonempty attachment.  Thus a single
proper interval gives a collar addition and no relative handle; two
intervals give one 1-handle.  Replacing every disc product by a 2-handle
would fail even this Euler check.

This theorem supplies local relative handles, not automatically an
admissible handle structure of uniform type.  It need not fix the entire
originally embedded core pointwise, because an attaching boundary disc
becomes interior during collar absorption.  Constructing the transported
pattern, arranging the required handle structure, and proving its
admissibility and global hierarchy bounds remain distinct tasks
\cite{Lackenby2026}.

\subsection{Implementation evidence and next steps}

The supplement contains a producer, an independent query constructor and
certificate checker, a finite oracle that actually walks polygon boundaries,
and a disc-replacement plan.  Source geometry is authenticated upstream by
\texttt{verify\_normal\_cut\_parallelity}.  The extra checker does not call
the producer's query builder, decoder or traversal.  Shared discovery caps
return no interface on exhaustion; verification caps cover all three
traces.  Cancellation propagates.

The saved audit passes 11 test methods with 184 literal geometric comparisons,
including 54 authenticated source comparisons.  It checks all 126 markings
of discs with one through six sides, annular and M\"obius charts, source
relabellings, rich 3-manifold fixtures, coincident endpoint addresses, and
forged cyclic, gap, fibre and weight data.  Exact JSON replay succeeds for
$2^{8192}$ parallel meridians in four tetrahedra with producer routines
disabled; the finite oracle refuses that input by its expansion cap.

Two complete example certificates accompany the audit.  The two-tetrahedron
source \texttt{fibonacci\_lst\_02}, with coordinate rows
$(0,0,1,1,1,0,0)$ and $(1,1,1,1,0,0,0)$, has one disc base supplied solely
by a selected 2-handle.  All eight ports lie consecutively around its
boundary, giving a full-circle 2-handle attachment.  Three parallel
meridians in the two-tetrahedron layered torus instead have two attached
disc bases with one proper core interval each: both are collar additions,
and two further ball products are unattached.  The prototype records these
local plans but performs no surgery.

The next concrete objectives are to certify the induced marked boundary
homeomorphisms; carry compressed pattern endpoints through these charts;
convert the relative additions to admissible uniform-type handles with a
proved size bound; and preserve coherent fibre data between different
boundary circles for non-disc bases.  The ordered interface removes an
explicit reconstruction gap, while these remaining obligations prevent
inferring a quasi-polynomial recognizer from this local operator alone.
